\documentclass{article}
\usepackage[pdftex]{graphicx}
\usepackage{xcolor}
\usepackage[letterpaper, margin = 2.5cm]{geometry}
\usepackage[T2A]{fontenc}
\usepackage[english, russian]{babel}
\usepackage{amsmath, amsfonts, amssymb}
\usepackage{hyperref, url}

\usepackage{minted}

\usepackage{parskip}

\newcommand{\solution}{\medskip\textbf{Решение}\medskip}

\begin{document}

\title{Домашнее задание 1}
\date{}
\author{SMR-2025 AIM, Никита Загайнов \href{https://t.me/V1adych}{@V1adych}}
\maketitle

\section*{Задача 1}

Покажите, что квадратная матрица $A$ коммутирует с $e^A$.

\solution

\begin{align*}
    A e^A & = A \left( I + \sum_{k=1}^{\infty} \frac{A^k}{k!} \right) \\
          & = A + \sum_{k=1}^{\infty} \frac{A^{k+1}}{k!}              \\
          & = \left( I + \sum_{k=1}^{\infty} \frac{A^k}{k!} \right) A \\
          & = e^A A
\end{align*}

\medskip

\section*{Задача 2}

\noindent Рассмотрим модель пружинного маятника на горизонтальной плоскости.

Его вектор состояния имеет вид

$x = \begin{pmatrix}
        p \\
        \dot{p}
    \end{pmatrix}$

Уравнение динамики системы в непрерывном случае таково:

$\dot{x} =  \begin{pmatrix}
        0             & 1 \\
        - \frac{k}{m} & 0
    \end{pmatrix} x$

Пусть симуляция этой системы производится как

$x_{k+1} = x_k + \Delta t \dot{x}_k$

Найдите, как изменяется полная энергия системы за один шаг симуляции.

\solution

Полная энергия системы
\[
    E = \frac{1}{2} m \dot{p}^2 + \frac{1}{2} k p^2
\]

Шаг симуляции

\[
    \begin{pmatrix}
        p_{k+1} \\
        \dot{p}_{k+1}
    \end{pmatrix}
    =
    \begin{pmatrix}
        p_k \\
        \dot{p}_k
    \end{pmatrix}
    + \Delta t \begin{pmatrix}
        0             & 1 \\
        - \frac{k}{m} & 0
    \end{pmatrix}
    \begin{pmatrix}
        p_k \\
        \dot{p}_k
    \end{pmatrix}
    =
    \begin{pmatrix}
        p_k + \Delta t \dot{p}_k \\
        \dot{p}_k - \frac{k}{m} \Delta t p_k
    \end{pmatrix}
\]

Изменение полной энергии системы за один шаг симуляции

\begin{align*}
    \Delta E & = E_{k+1} - E_k  =                                                                                    \\
             & = \frac{1}{2} m (\dot{p}_k - \frac{k}{m} \Delta t p_k)^2 + \frac{1}{2} k (p_k + \Delta t \dot{p}_k)^2 \\
             & - \frac{1}{2} m \dot{p}_k^2 - \frac{1}{2} k p_k^2 \\
             & = - \Delta t \dot{p}_k p_k + \frac{1}{2m} (k \Delta t p_k)^2 + k \Delta t \dot{p}_k p_k + \frac{k}{2} (\Delta t \dot{p}_k)^2
\end{align*}

\end{document}
